% Options for packages loaded elsewhere
% Options for packages loaded elsewhere
\PassOptionsToPackage{unicode}{hyperref}
\PassOptionsToPackage{hyphens}{url}
\PassOptionsToPackage{dvipsnames,svgnames,x11names}{xcolor}
%
\documentclass[
  ngerman,
  letterpaper,
  DIV=11]{scrreprt}
\usepackage{xcolor}
\usepackage{amsmath,amssymb}
\setcounter{secnumdepth}{3}
\usepackage{iftex}
\ifPDFTeX
  \usepackage[T1]{fontenc}
  \usepackage[utf8]{inputenc}
  \usepackage{textcomp} % provide euro and other symbols
\else % if luatex or xetex
  \usepackage{unicode-math} % this also loads fontspec
  \defaultfontfeatures{Scale=MatchLowercase}
  \defaultfontfeatures[\rmfamily]{Ligatures=TeX,Scale=1}
\fi
\usepackage{lmodern}
\ifPDFTeX\else
  % xetex/luatex font selection
\fi
% Use upquote if available, for straight quotes in verbatim environments
\IfFileExists{upquote.sty}{\usepackage{upquote}}{}
\IfFileExists{microtype.sty}{% use microtype if available
  \usepackage[]{microtype}
  \UseMicrotypeSet[protrusion]{basicmath} % disable protrusion for tt fonts
}{}
\makeatletter
\@ifundefined{KOMAClassName}{% if non-KOMA class
  \IfFileExists{parskip.sty}{%
    \usepackage{parskip}
  }{% else
    \setlength{\parindent}{0pt}
    \setlength{\parskip}{6pt plus 2pt minus 1pt}}
}{% if KOMA class
  \KOMAoptions{parskip=half}}
\makeatother
% Make \paragraph and \subparagraph free-standing
\makeatletter
\ifx\paragraph\undefined\else
  \let\oldparagraph\paragraph
  \renewcommand{\paragraph}{
    \@ifstar
      \xxxParagraphStar
      \xxxParagraphNoStar
  }
  \newcommand{\xxxParagraphStar}[1]{\oldparagraph*{#1}\mbox{}}
  \newcommand{\xxxParagraphNoStar}[1]{\oldparagraph{#1}\mbox{}}
\fi
\ifx\subparagraph\undefined\else
  \let\oldsubparagraph\subparagraph
  \renewcommand{\subparagraph}{
    \@ifstar
      \xxxSubParagraphStar
      \xxxSubParagraphNoStar
  }
  \newcommand{\xxxSubParagraphStar}[1]{\oldsubparagraph*{#1}\mbox{}}
  \newcommand{\xxxSubParagraphNoStar}[1]{\oldsubparagraph{#1}\mbox{}}
\fi
\makeatother


\usepackage{longtable,booktabs,array}
\newcounter{none} % for unnumbered tables
\usepackage{calc} % for calculating minipage widths
% Correct order of tables after \paragraph or \subparagraph
\usepackage{etoolbox}
\makeatletter
\patchcmd\longtable{\par}{\if@noskipsec\mbox{}\fi\par}{}{}
\makeatother
% Allow footnotes in longtable head/foot
\IfFileExists{footnotehyper.sty}{\usepackage{footnotehyper}}{\usepackage{footnote}}
\makesavenoteenv{longtable}
\usepackage{graphicx}
\makeatletter
\newsavebox\pandoc@box
\newcommand*\pandocbounded[1]{% scales image to fit in text height/width
  \sbox\pandoc@box{#1}%
  \Gscale@div\@tempa{\textheight}{\dimexpr\ht\pandoc@box+\dp\pandoc@box\relax}%
  \Gscale@div\@tempb{\linewidth}{\wd\pandoc@box}%
  \ifdim\@tempb\p@<\@tempa\p@\let\@tempa\@tempb\fi% select the smaller of both
  \ifdim\@tempa\p@<\p@\scalebox{\@tempa}{\usebox\pandoc@box}%
  \else\usebox{\pandoc@box}%
  \fi%
}
% Set default figure placement to htbp
\def\fps@figure{htbp}
\makeatother

\ifLuaTeX
  \usepackage{luacolor}
  \usepackage[soul]{lua-ul}
\else
  \usepackage{soul}
\fi

% definitions for citeproc citations
\NewDocumentCommand\citeproctext{}{}
\NewDocumentCommand\citeproc{mm}{%
  \begingroup\def\citeproctext{#2}\cite{#1}\endgroup}
\makeatletter
 % allow citations to break across lines
 \let\@cite@ofmt\@firstofone
 % avoid brackets around text for \cite:
 \def\@biblabel#1{}
 \def\@cite#1#2{{#1\if@tempswa , #2\fi}}
\makeatother
\newlength{\cslhangindent}
\setlength{\cslhangindent}{1.5em}
\newlength{\csllabelwidth}
\setlength{\csllabelwidth}{3em}
\newenvironment{CSLReferences}[2] % #1 hanging-indent, #2 entry-spacing
 {\begin{list}{}{%
  \setlength{\itemindent}{0pt}
  \setlength{\leftmargin}{0pt}
  \setlength{\parsep}{0pt}
  % turn on hanging indent if param 1 is 1
  \ifodd #1
   \setlength{\leftmargin}{\cslhangindent}
   \setlength{\itemindent}{-1\cslhangindent}
  \fi
  % set entry spacing
  \setlength{\itemsep}{#2\baselineskip}}}
 {\end{list}}
\usepackage{calc}
\newcommand{\CSLBlock}[1]{\hfill\break\parbox[t]{\linewidth}{\strut\ignorespaces#1\strut}}
\newcommand{\CSLLeftMargin}[1]{\parbox[t]{\csllabelwidth}{\strut#1\strut}}
\newcommand{\CSLRightInline}[1]{\parbox[t]{\linewidth - \csllabelwidth}{\strut#1\strut}}
\newcommand{\CSLIndent}[1]{\hspace{\cslhangindent}#1}

\ifLuaTeX
\usepackage[bidi=basic,shorthands=off]{babel}
\else
\usepackage[bidi=default,shorthands=off]{babel}
\fi
\ifLuaTeX
  \usepackage{selnolig} % disable illegal ligatures
\fi


\setlength{\emergencystretch}{3em} % prevent overfull lines

\providecommand{\tightlist}{%
  \setlength{\itemsep}{0pt}\setlength{\parskip}{0pt}}



 


\KOMAoption{captions}{tableheading}
\makeatletter
\@ifpackageloaded{tcolorbox}{}{\usepackage[skins,breakable]{tcolorbox}}
\@ifpackageloaded{fontawesome5}{}{\usepackage{fontawesome5}}
\definecolor{quarto-callout-color}{HTML}{909090}
\definecolor{quarto-callout-note-color}{HTML}{0758E5}
\definecolor{quarto-callout-important-color}{HTML}{CC1914}
\definecolor{quarto-callout-warning-color}{HTML}{EB9113}
\definecolor{quarto-callout-tip-color}{HTML}{00A047}
\definecolor{quarto-callout-caution-color}{HTML}{FC5300}
\definecolor{quarto-callout-color-frame}{HTML}{acacac}
\definecolor{quarto-callout-note-color-frame}{HTML}{4582ec}
\definecolor{quarto-callout-important-color-frame}{HTML}{d9534f}
\definecolor{quarto-callout-warning-color-frame}{HTML}{f0ad4e}
\definecolor{quarto-callout-tip-color-frame}{HTML}{02b875}
\definecolor{quarto-callout-caution-color-frame}{HTML}{fd7e14}
\makeatother
\makeatletter
\@ifpackageloaded{bookmark}{}{\usepackage{bookmark}}
\makeatother
\makeatletter
\@ifpackageloaded{caption}{}{\usepackage{caption}}
\AtBeginDocument{%
\ifdefined\contentsname
  \renewcommand*\contentsname{Inhaltsverzeichnis}
\else
  \newcommand\contentsname{Inhaltsverzeichnis}
\fi
\ifdefined\listfigurename
  \renewcommand*\listfigurename{Abbildungsverzeichnis}
\else
  \newcommand\listfigurename{Abbildungsverzeichnis}
\fi
\ifdefined\listtablename
  \renewcommand*\listtablename{Tabellenverzeichnis}
\else
  \newcommand\listtablename{Tabellenverzeichnis}
\fi
\ifdefined\figurename
  \renewcommand*\figurename{Abbildung}
\else
  \newcommand\figurename{Abbildung}
\fi
\ifdefined\tablename
  \renewcommand*\tablename{Tabelle}
\else
  \newcommand\tablename{Tabelle}
\fi
}
\@ifpackageloaded{float}{}{\usepackage{float}}
\floatstyle{ruled}
\@ifundefined{c@chapter}{\newfloat{codelisting}{h}{lop}}{\newfloat{codelisting}{h}{lop}[chapter]}
\floatname{codelisting}{Listing}
\newcommand*\listoflistings{\listof{codelisting}{Listingverzeichnis}}
\makeatother
\makeatletter
\makeatother
\makeatletter
\@ifpackageloaded{caption}{}{\usepackage{caption}}
\@ifpackageloaded{subcaption}{}{\usepackage{subcaption}}
\makeatother
\usepackage{bookmark}
\IfFileExists{xurl.sty}{\usepackage{xurl}}{} % add URL line breaks if available
\urlstyle{same}
\hypersetup{
  pdftitle={GO UNITE! Strategien zur Sammlung und Sichtbarkeit von Forschungsdaten der eigenen Institution},
  pdfauthor={Dorothea Strecker; Birte Cordes; Anna Daudrich; Jessica Hiller; Isabel Schuster; Martin Spenger},
  pdflang={de},
  colorlinks=true,
  linkcolor={blue},
  filecolor={Maroon},
  citecolor={Blue},
  urlcolor={Blue},
  pdfcreator={LaTeX via pandoc}}


\title{GO UNITE! Strategien zur Sammlung und Sichtbarkeit von
Forschungsdaten der eigenen Institution}
\author{Dorothea Strecker \and Birte Cordes \and Anna
Daudrich \and Jessica Hiller \and Isabel Schuster \and Martin Spenger}
\date{2026-10-01}
\begin{document}
\maketitle

\renewcommand*\contentsname{Inhaltsverzeichnis}
{
\setcounter{tocdepth}{2}
\tableofcontents
}

\bookmarksetup{startatroot}

\chapter*{Über dieses Dokument}\label{uxfcber-dieses-dokument}
\addcontentsline{toc}{chapter}{Über dieses Dokument}

\markboth{Über dieses Dokument}{Über dieses Dokument}

Dieses Dokument zeigt praxisnah, wie verstreut veröffentlichte
Forschungsdaten effizient identifiziert und der eigenen Einrichtung
zugeordnet werden können.

Der Teil \textbf{Strategien} enthält allgemeine Überlegungen zur
Konzeption und Durchführung des Monitorings von institutionellem
Forschungsdatenoutput.

Aktuell arbeiten wir ergänzend an \textbf{Notebooks}, die Code und Text
kombinieren und beispielhaft zeigen, wie bestimmte Datenquellen
abgefragt werden können. Die Notebooks werden voraussichtlich in Q1 2027
veröffentlicht.

Diese Version ist eine Aktualisierung und Überarbeitung von:

Putnings, M., Daudrich, A., Dille, N., Iarkaeva, A., Kühner, J., \&
Wollmann, A. (2024). GO UNITE! Empfehlungen zur Sammlung und
Sichtbarkeit von Forschungsdaten der eigenen Institution. Zenodo.
\url{https://doi.org/10.5281/zenodo.11102368}

\section*{Beitragende}\label{beitragende}
\addcontentsline{toc}{section}{Beitragende}

\markright{Beitragende}

Anastasiia Iarkaeva ;
\href{https://orcid.org/0000-0002-7370-1663}{0000-0002-7370-1663} ;
Berlin Institute of Health, Charité - Universitätsmedizin Berlin,
Deutschland

Markus Putnings ;
\href{https://orcid.org/0000-0002-6014-9048}{0000-0002-6014-9048} ;
Bibliothek, Technische Hochschule Deggendorf, Deutschland

Esther Asef ;
\href{https://orcid.org/0000-0003-2411-4953}{0000-0003-2411-4953} ;
Universitätsbibliothek, Charité - Universitätsmedizin Berlin,
Deutschland

Christian Malzer ;
\href{https://orcid.org/0009-0005-0522-7102}{0009-0005-0522-7102} ;
Bibliothek, Universität Würzburg, Deutschland

\part{Strategien}

\chapter{Einführung}\label{einfuxfchrung}

Dieses Dokument verfolgt das Ziel, praxisorientierte Strategien für die
Identifizierung, zentrale Erfassung und das Monitoring von
Forschungsdaten bereitzustellen, die von Angehörigen einer Institution
dezentral auf verschiedenen Plattformen -- beispielsweise Zenodo, Dryad
oder fachspezifischen Repositorien - publiziert wurden. Es baut auf dem
Dokument „GO UNITE! Empfehlungen zur Sammlung und Sichtbarkeit von
Forschungsdaten der eigenen Institution'' (Putnings u.~a. (2024)) auf,
wurde aber stark überarbeitet - einerseits aufgrund inhaltlicher
Weiterentwicklungen, andererseits um die Praxisrelevanz zu erhöhen.

In erster Linie richtet sich das Dokument an Personen, die mit dem
Monitoring und der Erfassung von Forschungsdaten betraut sind, etwa
Forschungsdatenbeauftragte, Mitarbeitende von
Forschungsinformationssystemen (FIS) sowie Bibliothekar:innen, aber auch
an alle, die sich sonst für dieses Thema interessieren.

In diesem Dokument umfasst der Begriff „Forschungsdaten'' sowohl Daten
im engeren Sinne -- etwa Mess- und Beobachtungsdaten, Laborergebnisse,
audiovisuelle Materialien, Texte, Umfragedaten, Sammlungsobjekte oder
Proben -- als auch den zugehörigen Code bzw. die Software, sofern diese
ebenfalls in Datenrepositorien veröffentlicht wurden. Der primäre Fokus
liegt auf Datensätzen, doch die hier beschriebenen Suchstrategien
differenzieren nicht zwischen Daten- und Codepublikationen. Für eine
gezielte und umfassende Erhebung veröffentlichter Forschungssoftware
müssten zusätzlich zu den hier genannten weitere, auf Software
spezialisierte Publikationsplattformen und Datenquellen berücksichtigt
werden.

Wie die Abfragen zum Zwecke des Monitorings dort funktionieren, ist
nicht Gegenstand dieses Dokuments.

\section{Motivation zur Durchführung des Datenmonitorings an der eigenen
Einrichtung}\label{motivation-zur-durchfuxfchrung-des-datenmonitorings-an-der-eigenen-einrichtung}

Ein Monitoring und begleitende Maßnahmen, wie z. B. eine Indexierung von
Forschungsdaten der eigenen Einrichtung, können maßgeblich dazu
beitragen, deren Auffindbarkeit und Nachnutzbarkeit im institutionellen
Kontext zu verbessern. Darüber hinaus erhöht sich dadurch die
Sichtbarkeit der publizierten Forschungsergebnisse über institutionelle
Grenzen hinaus, indem ihre Erschließung in externen Suchmaschinen,
Harvesting-Diensten, Forschungsinformationssystemen und weiteren
Aggregatoren erleichtert wird.

Die Sichtbarkeit von Forschungsdaten ist nicht nur im Hinblick auf deren
Nachnutzung von Bedeutung, sondern gewinnt auch vor dem Hintergrund der
zunehmenden Anerkennung von Forschungsdatenpublikationen als
eigenständige wissenschaftliche Leistung an Relevanz. Entsprechend
finden Forschungsdaten verstärkt Berücksichtigung in Verfahren der
Bewertung wissenschaftlicher Leistungen.

\section{Struktur des Dokuments}\label{struktur-des-dokuments}

Die Struktur des Dokuments orientiert sich am Prozess der
Identifizierung von Forschungsdaten. Entsprechend werden zunächst
notwendige Vorarbeiten und Voraussetzungen dargestellt, danach die
Durchführung beschrieben und schließlich weiterführende Perspektiven
sowie Entwicklungsmöglichkeiten aufgezeigt.

Ergänzend zu diesen Handlungsempfehlungen wird voraussichtlich in Q1
2027 ein separates Dokument veröffentlicht, das als Ergänzung zu diesem
Text zu verstehen ist. Darin wird anhand konkreter Praxisbeispiele
beschrieben, wie einzelne Plattformen nach den Forschungsdaten der
eigenen Institution durchsucht werden können.

Hilfreich sind darüber hinaus auch die Texte und Handreichungen, die von
weiteren Unterarbeitsgruppen der AG Sammlung und Sichtbarkeit von
Forschungsdaten der eigenen Institution entwickelt wurden.\footnote{\url{https://go-unite.de/index.php/aktivitaeten-und-publikationen/}
  (abgerufen am 30.09.2026)}

\chapter{Vorüberlegungen}\label{voruxfcberlegungen}

In diesem Kapitel werden Vorüberlegungen, die wissenschaftlichen
Einrichtungen bei der Konzeption von institutionellem
Forschungsdatenmonitoring helfen können, thematisiert. Grundsätzlich ist
es empfehlenswert, dass Einrichtungen, die ihren Forschungsdatenoutput
erfassen möchten, zunächst ihre Ziele festlegen und die vorhandenen
Ressourcen (insbesondere personelle Kapazitäten) bestimmen. Für die
Implementierung von Monitoring-Initiativen bei limitierten Ressourcen
werden Handlungsempfehlungen in einer Infobox vorgestellt.

\section{Ziele}\label{ziele}

Je nach institutionellen Rahmenbedingungen und strategischen
Anforderungen können Ziele und Anlässe eines Datenmonitorings variieren.
Das Monitoring kann ganz allgemein der Erfassung publizierter
Forschungsdaten dienen, etwa im Sinne einer Hochschulbibliografie, oder
darüber hinaus als Grundlage für strategische Entscheidungen und die
Weiterentwicklung des Forschungsdatenmanagements (FDM) genutzt werden.
Zu den möglichen Anlässen und Zielen eines Datenmonitorings zählen unter
anderem:

\begin{itemize}
\tightlist
\item
  Erweiterung von Dokumentationstypen in einem bereits bestehenden FIS
\item
  die Überprüfung der Umsetzung von Leitlinien
\item
  die Weiterentwicklung von Leitlinien
\item
  die Unterstützung des Berichtswesens und der strategischen
  Organisationsentwicklung
\item
  die Identifikation von Handlungsbedarfen im institutionellen FDM
\item
  die Erfolgskontrolle bereits umgesetzter FDM-Maßnahmen
\item
  die Berücksichtigung von Datenpublikationen in Einstellungs- und
  Evaluationsverfahren oder in der leistungsorientierten Mittelvergabe
\end{itemize}

Je nach Zielsetzung kann ein einmaliges oder ein permanentes Monitoring
erforderlich sein:

\begin{itemize}
\tightlist
\item
  Ein \textbf{einmaliges Monitoring} dient der anlassbezogenen Erfassung
  und Auswertung publizierter Forschungsdaten, beispielsweise zur
  Vorbereitung einer Evaluation, zur Entwicklung einer institutionellen
  FDM-Policy oder zur Analyse des aktuellen Stands des
  Forschungsdatenmanagements.
\item
  Im Gegensatz dazu verfolgt ein \textbf{permanentes Monitoring} das
  Ziel, publizierte Forschungsdaten kontinuierlich zu erfassen und
  auszuwerten, um Entwicklungen und Trends über einen längeren Zeitraum
  zu beobachten. Dadurch können beispielsweise die Wirksamkeit von
  FDM-Maßnahmen regelmäßig überprüft und neue Handlungsbedarfe
  frühzeitig erkannt werden. Um das Ziel einer möglichst vollständigen
  regelmäßigen Erfassung der FD-Publikationen zu unterstützen, können
  verschiedene Maßnahmen etabliert werden. Dazu zählen beispielsweise
  Empfehlungen oder Vorgaben zur Nutzung bestimmter Repositorien, die
  Erfassung von Metadaten in einem zentralen institutionellen
  Verzeichnis (z. B. FIS) sowie die Berücksichtigung der
  Forschungsdatenpublikation im Rahmen leistungsorientierter
  Mittelvergabe.
\end{itemize}

\section{Herausforderungen}\label{herausforderungen}

Die Erfassung von Forschungsdatenpublikationen einer Einrichtung ist mit
technischen und institutionellen Herausforderungen verbunden.
Nachfolgend werden die wichtigsten Hürden dargestellt, um sie in der
Konzeption zu berücksichtigen und realistisch planen zu können.

\subsection{Technische
Herausforderungen}\label{technische-herausforderungen}

Zunächst werden nun einige technische Herausforderungen vorgestellt,
welche die Erfassung des Forschungsdatenoutputs beeinträchtigen. Wie
diese Herausforderungen in der Praxis adressiert werden können, wird in
\href{./3_Durchführung.qmd}{Kapitel 3} gezeigt.

\begin{enumerate}
\def\labelenumi{\arabic{enumi}.}
\tightlist
\item
  Verschiedene Datenquellen und Suchstrategien
\end{enumerate}

\begin{itemize}
\tightlist
\item
  \textbf{Auswahl geeigneter Datenquellen}: Welche Objekte deckt die
  Datenquelle ab? Enthält die Datenquelle die Metadaten, die für den
  Zweck erforderlich sind? Bietet die Datenquelle geeignete Zugänge und
  Ausgabeformate an? In \hyperref[suchstrategien]{Kapitel 3.2} werden
  einige zentrale Datenquellen vorgestellt. Grundsätzlich sollten
  Datenquellen vor der Auswahl gründlich geprüft werden. Dazu zählt
  neben der Konsultation der Dokumentation auch ein Testdurchlauf, um
  Zugangswege und Ausgabeformate zu prüfen.
\item
  \textbf{Auswahl geeigneter Suchstrategien}: Gibt es bereits
  Anhaltspunkte darüber, wo und wie Forschende der Einrichtung ihre
  Daten publizieren (beispielsweise aus Bedarfserhebungen)? Welche
  Suchstrategien können mit den gewählten Datenquellen und vorhandenen
  Ressourcen umgesetzt werden? In
  \hyperref[grundprinzipien-und-schematischer-ablauf]{Kapitel 3.1}
  werden verschiedene Suchstrategien vorgestellt. Empfehlenswert ist
  nach Möglichkeit die Kombination verschiedener Suchstrategien.
\end{itemize}

\begin{enumerate}
\def\labelenumi{\arabic{enumi}.}
\setcounter{enumi}{1}
\tightlist
\item
  Eingeschränkte technische Kapazitäten
\end{enumerate}

\begin{itemize}
\tightlist
\item
  \textbf{Auswahl geeigneter Werkzeuge}: Sind programmatische Abfragen
  möglich? Welche Werkzeuge stehen für die Weiterverarbeitung der Daten
  zur Verfügung (z. B. bestimmte Skriptsprachen oder Text- und
  Tabellenbearbeitungsprogramme)?
\item
  \textbf{Bestimmen des Ressourceneinsatzes}: Welche personellen
  Kapazitäten stehen zur Verfügung -- wie viele Personen können sich mit
  dem Thema beschäftigen und wie viel Arbeitszeit können sie darauf
  verwenden?
\end{itemize}

\begin{enumerate}
\def\labelenumi{\arabic{enumi}.}
\setcounter{enumi}{2}
\tightlist
\item
  Vorgaben des Zielsystems
\end{enumerate}

\begin{itemize}
\tightlist
\item
  \textbf{Inputformate und Metadatenfelder}: Welche Anforderungen stellt
  das Zielsystem?
\item
  \textbf{Importmöglichkeiten}: Verfügt das Zielsystem über
  Schnittstellen und wenn ja, welche? Welcher Maß an Automatisierung
  kann ereicht werden und welcher manueller Aufwand ist weiterhin
  notwendig?
\end{itemize}

\subsection{Institutionelle
Herausforderungen}\label{institutionelle-herausforderungen}

Arbeitsgruppe 1 (Scoping the Needs of Open Science Monitoring) der Open
Science Monitoring Initiative hat festgestellt, dass institutionelle
Barrieren die größten Herausforderungen bei der Umsetzung von
Monitoringaktivitäten darstellen (Salisu und Cortés, 2026). Dazu zählen
begrenzte Ressourcen sowie mangelnde Governance und Nachhaltigkeit.
Darum ist es wichtig, ausreichend personelle Kapazitäten für das
Vorhaben einzuplanen und einen realistischen Zeitrahmen zu setzen. Zum
Einstieg können die folgenden Ansätze in Betracht gezogen werden:

\begin{enumerate}
\def\labelenumi{\arabic{enumi}.}
\tightlist
\item
  Wenn die verfügbaren Ressourcen beschränkt oder unklar sind, kann
  zunächst stufenweise vorgegangen werden. Dabei wird zuerst eine
  Minimalumsetzung anvisiert, beispielsweise in Form der Suchstrategie
  „direkte Suche'' (s. \hyperref[direkte-suche]{Kapitel 3.2.1}) mit
  lediglich einer Datenquelle. Dieser Ansatz kann ggf. später erweitert
  werden, etwa indem mehr Datenquellen einbezogen werden, die
  Suchstrategie modifiziert oder die Erfassung zusätzlich auf
  Forschungssoftware ausgeweitet wird.
\item
  Empfehlenswert ist zunächst eine einmalige Durchführung mit
  detaillierter Dokumentation der Arbeitsschritte. Die so entstehende
  Basis unterstützt eine spätere Verstetigung und zeigt, ob
  institutionelles Forschungsdatenmonitoring grundsätzlich umsetzbar
  ist, bevor ein kontinuierlicher Betrieb etabliert wird. Strategische
  Fragen wie die Ansiedlung der Zuständigkeit oder eine weitgehende
  Automatisierung können vorerst zurückgestellt und erst bei einer
  Verstetigung angegangen werden.
\end{enumerate}

\begin{tcolorbox}[enhanced jigsaw, arc=.35mm, bottomrule=.15mm, bottomtitle=1mm, breakable, colback=white, colbacktitle=quarto-callout-tip-color!10!white, colframe=quarto-callout-tip-color-frame, coltitle=black, left=2mm, leftrule=.75mm, opacityback=0, opacitybacktitle=0.6, rightrule=.15mm, title=\textcolor{quarto-callout-tip-color}{\faLightbulb}\hspace{0.5em}{Infobox: Welche Optionen existieren bei beschränkten oder unklaren
Ressourcen?}, titlerule=0mm, toprule=.15mm, toptitle=1mm]

„Perfect is the enemy of good.'' - Natürlich ist eine vollständige
Erfassung wünschenswert. Aus den oben beschriebenen Gründen ist sie
jedoch nicht immer realistisch. Wenn die „perfekte'' Lösung nicht
möglich ist, können die folgenden Strategien ggf. dabei helfen, einen
ersten Einstieg in das Monitoring von Forschungsdatenpublikationen zu
erreichen:

\begin{itemize}
\tightlist
\item
  \textbf{Ressourcen mobilisieren}: Gibt es Personen oder Abteilungen an
  der Einrichtung, die ebenfalls Interesse an der Erfassung von
  Forschungsdatenpublikationen haben (Data Stewards, Fachreferent:innen,
  Entwicklungsplanung o. ä.) und die ggf. Ressourcen beisteuern können?
\item
  \textbf{Auf eine oder wenige Datenquellen beschränken}:
  Steckbriefartige Übersichten über ausgewählte Datenquellen sind am
  Ende von \href{./3_Durchführung.qmd}{Kapitel 3} zu finden. Relevant
  für die Entscheidung können u. a. die folgenden Aspekte sein:

  \begin{itemize}
  \tightlist
  \item
    \textbf{Scope} \footnotemark{}: Woher bezieht die Datenquelle ihre
    Daten? Welche weiteren Datenquellen sind inkludiert? Gibt es
    fachspezifische Quellen oder bestimmte Projekte, die auf jeden Fall
    enthalten sein sollen? Wenn möglichst viele Datenpublikationen
    gefunden werden sollen, empfiehlt sich ein Aggregator wie DataCite.

    \begin{itemize}
    \tightlist
    \item
      \textbf{Abfragemöglichkeiten}: Passen sie zu den vorhandenen
      technischen Möglichkeiten? Kann z. B. die geplante Suchanfrage
      über die Nutzeroberfläche gestellt werden, falls keine Nutzung der
      API möglich ist? Ist ein Komplettexport der Trefferliste in einem
      geeigneten Dateiformat möglich?
    \item
      \textbf{Metadaten}: Sind bestimmte Metadaten besonders wichtig (z.
      B. Affiliations-Metadatenfeld)? Können diese bei der gewählten
      Datenquelle abgefragt werden?
    \item
      \textbf{Offenheit}: Ist die Datenquelle frei von Zugangs- und
      Nutzungsbeschränkungen?
    \end{itemize}
  \end{itemize}
\item
  \textbf{Suche einschränken}: Lässt sich die Suche einschränken,
  beispielsweise auf Forschungsdatenpublikationen mit DOI oder auf
  bestimmte Metadatenfelder?
\item
  \textbf{Suche präzisieren}: Wenige, konkrete Zielfragen formulieren
  und die Datensammlung und -auswertung darauf ausrichten
\item
  \textbf{Vorhandene Verzeichnisse berücksichtigen}: Können bestehende
  Verzeichnisse (z. B. FIS, Hochschulbibliografie, Repositorien) bei dem
  Vorhaben helfen?
\item
  \textbf{Vorgehen dokumentieren}: Durch eine detaillierte Dokumentation
  des Vorgehens kann auch angesichts schlecht planbarer Zeiträume sowie
  personeller und finanzieller Engpässe später immer wieder auf dem
  bisherigen Fortschritt aufgebaut werden
\end{itemize}

\end{tcolorbox}

\footnotetext{Mit „Scope'' ist hier „Geltungsbereich'' gemeint. Da der
englische Begriff im Deutschen geläufig ist, verwenden wir ihn im
gesamten Dokument.}

\chapter{Durchführung}\label{durchfuxfchrung}

Das folgende Kapitel behandelt die praktische Durchführung des
Monitorings des institutionellen Forschungsdatenoutputs. Es stellt
leitende Grundprinzipien vor und fasst den schematischen Ablauf von
Monitoringaktivitäten zusammen. Zudem werden zwei Suchstrategien - die
direkte und die indirekte Suche - beschrieben. Abschließend werden
exemplarische Datenquellen für die direkte Suche präsentiert.

\section{Grundprinzipien und schematischer
Ablauf}\label{grundprinzipien-und-schematischer-ablauf}

Monitoringaktivitäten sollten grundsätzlich individuell auf die
jeweilige Einrichtung zugeschnitten werden, einem klar definierten
Anlass folgen, bestehende Rahmenbedingungen - etwa Leitlinien und
Fächerprofile - berücksichtigen und alle relevanten Akteure einbeziehen,
wobei diese je nach Einrichtung variieren.

Einige Grundprinzipien sind auf alle Monitoringaktivitäten im Kontext
Open Science anwendbar. 2025 veröffentlichte eine Gruppe internationaler
Expert:innen die ``Principles of Open Science Monitoring'' (Bobrov u.~a.
(2025)). Dieses Dokument formuliert Anforderungen an
Monitoringaktivitäten wie die Offenheit und Qualität der genutzten
Datenquellen oder die Berücksichtigung spezifischer Fächerkulturen.
Einrichtungen, die ihren Forschungsdatenoutput messen möchten, wird
nahegelegt, ihre Aktivitäten in regelmäßigen Abständen mit den
Principles of Open Science Monitoring abzugleichen.

Die praktische Ausgestaltung des Monitorings hängt neben den Zielen von
weiteren Faktoren ab, etwa den vorhandenen Ressourcen. In den meisten
Fällen wird der Ablauf die folgenden Schritte beinhalten:

\begin{enumerate}
\def\labelenumi{\arabic{enumi}.}
\tightlist
\item
  \textbf{Die Auswahl geeigneter Suchstrategien}. Diese werden in
  \hyperref[suchstrategien]{Kapitel 3.2} näher vorgestellt.
\item
  \textbf{Die Auswahl geeigneter Datenquellen}. Einige Datenquellen
  werden exemplarisch in
  \hyperref[datenquellen-fuxfcr-die-direkte-suche]{Kapitel 3.3}
  präsentiert.
\item
  \textbf{Die Durchführung der Abfragen}. Um die Überprüfbarkeit der
  Ergebnisse und die Reproduzierbarkeit des Verfahrens sicherzustellen,
  sollte die Durchführung gründlich dokumentiert werden, etwa indem
  Skripte für die Abfragen erstellt und abgespeichert werden.
  Insbesondere die verwendeten Suchbegriffe, Filter und das Abfragedatum
  sind zu dokumentieren, etwa für die Einschränkung auf bestimmte
  Publikationstypen oder Namensvarianten der Einrichtung. Einige
  Beispiele werden ergänzend zu diesem Dokument voraussichtlich ab Q1
  2027 in Notebooks zur Verfügung gestellt.
\item
  \textbf{Die Bereinigung der Ergebnisse}. Nach der Durchführung der
  Abfragen muss die Treffermenge bereinigt werden. In erster Linie geht
  es hierbei um das Entfernen von Dubletten, die Zusammenfassung von
  Versionen und gegebenenfalls die einheitliche Ansetzung von Werten in
  Freitextfeldern. Die Vereinheitlichung von Werten in Freitextfeldern,
  etwa für die Namen von Autor:innen, Repositorien und
  Förderorganisationen, kann sehr aufwendig sein (Kirsch und Wink
  (2026)). Wenn Forschungsdaten über die Suche in Volltexten ermittelt
  werden, zählt dazu auch der Ausschluss von Fällen, die lediglich
  ``Datennachnutzung'' darstellen.
\item
  \textbf{Die Bewertung der Ergebnisse}. Die Treffermenge sollte auf
  Korrektheit überprüft werden, um falschpositive Treffer
  auszuschließen. Dazu zählt insbesondere eine Prüfung der
  Affiliationsangaben, um sicherzustellen, dass die gefundenen
  Forschungsdaten tatsächlich der Einrichtung zugeordnet werden können.
\end{enumerate}

Nach diesen Prozessschritten können die Ergebnisse schließlich für die
Weiterverarbeitung im Zielsystem aufbereitet werden (s.
\href{./4_Verwendung.qmd}{Kapitel 4}).

Welche weiteren Maßnahmen zur Erleichterung zukünftiger
Monitoringaktivitäten ergriffen werden können, wird in
\href{./5_Vorbereitung}{Kapitel 5} genauer dargelegt.

\section{Suchstrategien}\label{suchstrategien}

Die grundlegende Herausforderung beim Monitoring des
Forschungsdatenoutputs einer Einrichtung ist die gesicherte Herstellung
eines Bezugs zwischen Forschungsdatensätzen und der Einrichtung. Diese
Zuordnung kann durch zwei verschiedene Ansätze realisiert werden. Durch
eine Kombination der beiden Ansätze erhöht sich die Treffermenge.

Der Bezug zwischen einer Einrichtung und einem Datensatz lässt sich
\textbf{direkt} herstellen, indem Metadatensätze, die Forschungsdaten
beschreiben, durchsucht werden -- etwa nach Angaben zu Personen, welche
die Datensätze erstellt haben, und ihrer Affiliation. Gee (2026)
beschreibt die Durchführung dieser Suchstrategie umfassend am Beispiel
der University of Texas. Häufig fehlen diese Angaben jedoch in den
Metadaten, beispielsweise weil sie bei der Übergabe an ein Repositorium
nicht mit überliefert werden (Robinson-Garcia u.~a. (2017)).

Darum ist ergänzend die \textbf{indirekte} Suche über andere
Forschungsoutputs, die der Einrichtung zugeordnet werden können,
sinnvoll. Dazu zählen beispielsweise Zeitschriftenartikel, die auf
zugrundeliegende Datensätze verweisen. Der Bezug zwischen einer
Einrichtung und einem Datensatz wird dann indirekt über diesen
Forschungsoutput hergestellt. Cohen u.~a. (2026) verdeutlichen am
Beispiel der Charité Universitätsmedizin Berlin, wie mit diesem Ansatz
Forschungsdatenpublikationen identifiziert werden können.

Im Folgenden werden die beiden Ansätze näher beschrieben.

\subsection{Direkte Suche}\label{direkte-suche}

Forschungsdatensätze können einer Einrichtung direkt zugeordnet werden,
wenn die Metadaten, die den Datensatz beschreiben, entsprechende Angaben
enthalten. Für diese Suchstrategie können die unten beschriebenen
Datenquellen genutzt werden.

Grundsätzlich sind hier zwei Ansätze denkbar: die Suche in
Affiliationsangaben oder die Suche nach Autor:innen. Einige Datenquellen
nutzen in ihren Metadaten persistente Identifikatoren für Institutionen
(z. B. ROR IDs) und Autor:innen (z. B. ORCID IDs). Diese ermöglichen
eine eindeutige Zuordnung zur Einrichtung bzw. Person, darum wäre diese
Option besonders wünschenswert für das Monitoring. Da persistente
Identifikatoren allerdings nicht immer mit angegeben werden, ist
ergänzend eine Suche in Freitextfeldern zu empfehlen, etwa im
Freitextfeld für den Namen der Institution. Bei der Suche in
Affiliationsangaben sollten möglichst viele verschiedene
Ansetzungsformen und Schreibweisen des Einrichtungsnamens genutzt
werden, zum Beispiel neben ``Humboldt-Universität zu Berlin'' auch ``HU
Berlin'' oder ``Humboldt University Berlin''. Namensvarianten finden
sich beispielsweise im ROR-Eintrag einer Einrichtung oder auch in der
institutionellen Affiliationsrichtlinie, falls vorhanden.\footnote{Beispiel
  für Namensvarianten im ROR-Profil der Humboldt-Universität zu Berlin:
  \url{https://ror.org/01hcx6992}}

Jede Suche wird verschiedene Treffermengen erzeugen, die im Anschluss
dedupliziert werden sollten.

Bei Datenquellen, die nicht auf Forschungsdaten spezialisiert sind,
sollte die Suche auf Publikationstypen eingeschränkt werden, die
Gegenstand des Monitorings sind. Nach welchen Werten zu suchen ist,
hängt davon ab, wie die Datenquelle Publikationstypen beschreibt. Wenn
kein kontrolliertes Vokabular genutzt wird, ist auch hier
empfehlenswert, nach mehreren Varianten und Schreibweisen zu suchen,
beispielsweise nach ``data'', ``dataset'' und ``data set''. Bei der Wahl
von Suchbegriffen kann z. B. das Vokabular COAR Resource Types zurate
gezogen werden.\footnote{\url{https://vocabularies.coar-repositories.org/resource_types/}
  (abgerufen am 30.09.2026)}

Die größte Limitation dieses Ansatzes ist, dass er wesentlich von der
Metadatenqualität der genutzten Datenquellen abhängt (Robinson-Garcia
u.~a. (2017)). Fehlende oder fehlerhafte Werte können das Ergebnis
verfälschen.

\subsection{Indirekte Suche}\label{indirekte-suche}

Da Affiliationsangaben in Metadaten, die Forschungsdaten beschreiben,
häufig fehlen (Habermann (2025)), kann als weiterer Ansatz eine
indirekte Suche über Verweise auf Datensätze in Textpublikationen wie
Zeitschriftenartikel oder Data Paper herangezogen werden. Die
Verknüpfung zwischen einem Datensatz und einer Einrichtung erfolgt bei
diesem Ansatz mittelbar über eine Textpublikation, die sowohl einer
Einrichtung zugeordnet werden kann als auch auf einen Datensatz
verweist, der die vorgestellten Ergebnisse stützt. Einige Einrichtungen
sammeln bereits den Publikationsoutput der Einrichtungsangehörigen, z.
B. in einem FIS oder in Form einer Hochschulbibliografie. Diese
Sammlungen können als Basis dienen, zugehörige Datensätze zu finden.
Auch der OA Monitor\footnote{\url{https://open-access-monitor.de/}
  (abgerufen am 30.09.2026)} kann als Ausgangspunkt für diese
Suchstrategie genutzt werden, da diese Datenquelle Textpublikationen mit
wissenschaftlichen Einrichtungen verknüpft.

Der Verweis auf einen Datensatz kann sowohl im Literaturverzeichnis der
Textpublikation als auch an anderer Stelle im Volltext, zum Beispiel im
Methodenteil oder in einem Data Availability Statement, erfolgen.
Gregory u.~a. (2023) (S. 625) unterscheiden zwischen:

\begin{itemize}
\tightlist
\item
  data citation: Verweis auf Datensätze im Literaturverzeichnis
\item
  data mention: Verweis auf Datensätze an anderer Stelle
\item
  indirect data citation: Verweis auf eine Textpublikation statt des
  Datensatzes im Literaturverzeichnis
\end{itemize}

``Data citations'' werden von verschiedenen Datenquellen erfasst, zum
Beispiel vom Data Citation Corpus (DataCite und Make Data Count (2025))
oder von DataCite (DataCite (2026)). Allerdings gibt es große
Abweichungen zwischen den Datenquellen in der Abdeckung von
Zitationsbeziehungen (Gerasimov u.~a. (2024) ; Strecker u.~a. (2025)).
Darum sollten unbedingt mehrere Datenquellen kombiniert werden, um
Lücken auszugleichen.

Einige Einrichtungen extrahieren darüber hinaus bereits ``data
mentions'', also Verweise auf Datensätze in Textpublikationen, um den
Forschungsdatenoutput zu bestimmen. Hier ist insbesondere ein Projekt am
QUEST Center for Responsible Research am Berlin Institute of Health at
Charité zu nennen, das einen eigenen Workflow und Algorithmus für diesen
Ansatz entwickelt hat (Iarkaeva u.~a. (2024) ; Riedel und Nachev
(2026)). Die Methode ist jedoch für die Charité Universitätsmedizin
Berlin und mit Fokus auf die Medizin entwickelt worden. Folglich kann
diese Vorgehensweise nicht ohne Anpassung auf andere Institutionen oder
Fachbereiche übertragen werden. Diese Anpassung hat das Team am Beispiel
der Geowisseschaften für die Fachbibliothek Geowissenschaften an der
Freien Universität Berlin\footnote{\url{https://www.geo.fu-berlin.de/bibliotheken/geo/index.html}
  (abgerufen am 30.09.2026)} demonstriert (Iarkaeva u.~a. (o.~J.)). Auch
an der Universitätsbibliothek Mannheim wurden Workflows und Software
entwickelt, um über die Extraktion von Verweisen auf Datensätzen aus
Publikationen in der Universitätsbibliografie den Forschungsdatenoutput
zu erfassen (Universitätsbibliothek Mannheim (2026)).

Im Rahmen der Prüfung der Treffermenge bei der Suchstrategie „indirekte
Suche'' ist sicherzustellen, dass Fälle von ``Datennachnutzung''
ausgeschlossen werden sollten. Leider ist der Begriff
„Datennachnutzung'' nicht eindeutig definiert (Sandt u.~a. (2019)).
Näherungsweise können Datensätze, die nicht von den Autor:innen der
Textpublikation selbst erstellt wurden (also ``nachgenutzt'' wurden),
herausgefiltert werden. Das ist wichtig, da sie nicht zum
Forschungsoutput der Einrichtung gehören und das Ergebnis verfälschen
würden. Wie das Herausfiltern nachgenutzter Datensätze durch einen
Abgleich der Autor:innen erfolgen kann, beschreiben Cohen u.~a. (2026)
für die Charité Universitätsmedizin Berlin.

Dieser Ansatz bringt allerdings einige Limitationen mit sich. Zunächst
hängt er wesentlich von den Zitationspraktiken der Forschenden ab.
Studien zeigen, dass Forschende Datensätze oft nicht oder nur indirekt
zitieren (Gregory u.~a. (2023)). Auch hängt er von der Verfügbarkeit von
Volltexten ab: Textpublikationen können nur nach Verweisen auf
Datensätze durchsucht werden, wenn Zugriff auf die Volltexte besteht.
Darum ist es möglich, dass dieser Ansatz für bestimmte Einrichtungen
oder Fächerprofile nur unvollständige Ergebnisse liefert. Insgesamt ist
die indirekte Suche ressourcenintensiver als die direkte Suche (Gee
(2026)).

\section{Datenquellen für die direkte
Suche}\label{datenquellen-fuxfcr-die-direkte-suche}

Es gibt verschiedene Arten von Datenquellen, die für das Monitoring über
die direkte Suche genutzt werden können. Datenquellen können
beispielsweise unterschieden werden in individuelle Repositorien und
Aggregatoren, die Metadaten aus mehreren Quellen zusammenführen. Einige
Datenquellen sind spezialisiert auf Forschungsdaten, andere weisen
darüber hinaus auch weitere Forschungsoutputs nach.

Grundsätzlich empfiehlt es sich, zunächst die für die eigene Einrichtung
relevanten Datenquellen zu identifizieren. Dabei sollten Abdeckung,
Metadatenschema, Zugänge und Ausgabeformate berücksichtigt werden. Um
die Transparenz der Ergebnisse sicherzustellen, sollten nach Möglichkeit
offene Datenquellen genutzt werden, wie die UNESCO Recommendations on
Open Science nahelegen: ``The monitoring of open science should be
explicitly kept under public oversight, including the scientific
community, and whenever possible supported by open non-proprietary and
transparent infrastructures.'' (\emph{{UNESCO} Recommendation on Open
Science} (2021) ; S. 34)

Es ist empfehlenswert, bei der Suche zunächst auf Aggregatoren zu
setzen, da dies Ressourcen spart und einheitlichere Treffermengen
liefert. Je nach Bedarf sollte darüber hinaus auch in individuellen
Repositorien recherchiert werden -- v.a. dann, wenn bekannt ist, dass
Einrichtungsangehörige auch kleinere oder disziplinspezifische
Repositorien nutzen, die bspw. keine DOIs vergeben (etwa
Sequenzdatenbanken). Darüber hinaus kann eine Direktsuche in
Repositorien zur Evaluation der Vollständigkeit der über Suchmaschinen
gefundenen Publikationen herangezogen werden.

Da sich Datenquellen im Laufe der Zeit verändern können, etwa wenn das
Metadatenschema aktualisiert wird, wird an dieser Stelle auf Beispiele
für spezifische Abfragen verzichtet. Stattdessen werden wir
voraussichtlich in Q1 2027 interaktive Notebooks veröffentlichen, die
angepasst und ergänzt werden können.

Im Folgenden werden einige zentrale offene Datenbanken vorgestellt, die
als Aggregatoren einen guten Einstieg in das Monitoring bieten können.
Kriterien für die Auswahl von Datenquellen sind in der Infobox in
\hyperref[voruxfcberlegungen]{Kapitel 2} zu finden.

\subsection{BASE}\label{base}

{\def\LTcaptype{none} % do not increment counter
\begin{longtable}[]{@{}
  >{\raggedright\arraybackslash}p{(\linewidth - 2\tabcolsep) * \real{0.1327}}
  >{\raggedright\arraybackslash}p{(\linewidth - 2\tabcolsep) * \real{0.8673}}@{}}
\toprule\noalign{}
\begin{minipage}[b]{\linewidth}\raggedright
Typ
\end{minipage} & \begin{minipage}[b]{\linewidth}\raggedright
Suchmaschine ; Aggregator
\end{minipage} \\
\midrule\noalign{}
\endhead
\bottomrule\noalign{}
\endlastfoot
Beschreibung & \url{https://www.base-search.net/about/en/index.php} \\
Abdeckung & Metadaten zu Forschungsoutputs, die von Datengebenden über
OAI-PMH-Schnittstellen zum Harvesting angeboten werden;
\href{https://www.base-search.net/about/de/about_sources_date.php}{Datenlieferanten} \\
Metadatenschema & \href{https://oai.base-search.net/\#oai-dc}{oai-dc} ;
\href{https://oai.base-search.net/\#base-dc}{base-dc} \\
Zugänge & GUI mit
``\href{https://www.base-search.net/Search/Advanced}{Erweiterter
Suche}'' ; API (\href{https://oai.base-search.net/}{OAI-PMH}) ; Der
Zugang via API ist ausschließlich für nicht-kommerzielle Zwecke
gestattet und erfordert einen API-Key, der nach erfolgter Registrierung
zur Verfügung gestellt wird (siehe
\url{https://ubib-baseapi.ub.uni-bielefeld.de/}, zuletzt aufgerufen am
02.09.2026). Das Anfragenlimit beträgt eine Anfrage pro Sekunde. \\
Suchfelder & Publikationstyp kann auf Forschungsdaten eingeschränkt
werden (doctype:7) ;Kein Metadatenfeld für Affiliation / ROR ID
vorhanden, daher muss mit Freitext gesucht werden. \\
Ausgabeformate & GUI: verschiedene Formate, u. a. JSON (kein csv) ; API:
XML \\
\end{longtable}
}

\subsection{DataCite}\label{datacite}

{\def\LTcaptype{none} % do not increment counter
\begin{longtable}[]{@{}
  >{\raggedright\arraybackslash}p{(\linewidth - 2\tabcolsep) * \real{0.1327}}
  >{\raggedright\arraybackslash}p{(\linewidth - 2\tabcolsep) * \real{0.8673}}@{}}
\toprule\noalign{}
\begin{minipage}[b]{\linewidth}\raggedright
Typ
\end{minipage} & \begin{minipage}[b]{\linewidth}\raggedright
DOI registration agency ; Aggregator
\end{minipage} \\
\midrule\noalign{}
\endhead
\bottomrule\noalign{}
\endlastfoot
Beschreibung & \url{https://datacite.org/mission/} \\
Abdeckung & Metadaten zu Forschungsoutputs, für die eine DOI bei
DataCite vergeben wurde, sowie verwandte Publikationen
(\href{https://support.datacite.org/docs/works-in-datacite-commons}{Quelle}) \\
Metadatenschema & \href{https://doi.org/10.14454/qdd3-ps68}{DataCite
Metadata Schema v4.7} \\
Zugänge & DataCite Commons (\href{https://commons.datacite.org/}{GUI}) ;
DataCite \href{https://support.datacite.org/docs/api}{REST API} mit
\href{https://support.datacite.org/docs/rate-limit}{rate limits} ;
DataCite metadata dashboards
(\href{https://metadata.datacite.org/}{GUI}) ; Die öffentliche REST API
kann ohne Authentifizierung genutzt werden und wird von DataCite für das
Abrufen und Suchen von DOI-Metadaten empfohlen. \\
Suchfelder & Im Tab ``Organizations'' kann mit ROR ID gesucht werden,
dies führt zum Dashboard für die Institution. Die Treffermenge kann dann
auf den work type „Dataset`` eingeschränkt werden. \\
Ausgabeformate & DataCite Commons: Download einer csv-Datei; DataCite
REST API: JSON \\
\end{longtable}
}

\subsection{OpenAIRE}\label{openaire}

{\def\LTcaptype{none} % do not increment counter
\begin{longtable}[]{@{}
  >{\raggedright\arraybackslash}p{(\linewidth - 2\tabcolsep) * \real{0.1327}}
  >{\raggedright\arraybackslash}p{(\linewidth - 2\tabcolsep) * \real{0.8673}}@{}}
\toprule\noalign{}
\begin{minipage}[b]{\linewidth}\raggedright
Typ
\end{minipage} & \begin{minipage}[b]{\linewidth}\raggedright
Aggregator
\end{minipage} \\
\midrule\noalign{}
\endhead
\bottomrule\noalign{}
\endlastfoot
Beschreibung & \url{https://www.openaire.eu/faqs} \\
Abdeckung & Zahlreiche europäische Datenquellen, darunter sowohl
DOI-Agenturen und Aggregatoren als auch Repositorien
(\href{https://explore.openaire.eu/search/find/dataproviders}{Quellenliste}) \\
Metadatenschema &
\href{https://doi.org/10.5281/zenodo.4238938}{JSON-Schema} \\
Zugänge & OpenAIRE Explore \href{https://explore.openaire.eu/}{(GUI)} ;
OpenAIRE Graph
(\href{https://graph.openaire.eu/docs/apis/graph-api/}{API}) ; Die
öffentliche Graph API kann sowohl mit als auch ohne Authentifizierung
genutzt werden. Das Anfragenlimit beträgt bei authentifiziertem Zugang
7200 Anfragen pro Stunde, bei unauthentifiziertem Zugang 60 Anfragen pro
Stunde. Zur Authentifizierung ist ein API-Key erforderlich; dieser wird
nach Erstellung eines Nutzeraccounts zur Verfügung gestellt. \\
Suchfelder & Institutionen sind als kontrolliertes Vokabular angelegt
und können in der Advanced Search in Research products mit
``Organization is'' ausgewählt werden. Filter auf ``Research Data''
möglich \\
Ausgabeformate & GUI: csv \\
\end{longtable}
}

\subsection{OpenAlex}\label{openalex}

{\def\LTcaptype{none} % do not increment counter
\begin{longtable}[]{@{}
  >{\raggedright\arraybackslash}p{(\linewidth - 2\tabcolsep) * \real{0.1327}}
  >{\raggedright\arraybackslash}p{(\linewidth - 2\tabcolsep) * \real{0.8673}}@{}}
\toprule\noalign{}
\begin{minipage}[b]{\linewidth}\raggedright
Typ
\end{minipage} & \begin{minipage}[b]{\linewidth}\raggedright
Aggregator
\end{minipage} \\
\midrule\noalign{}
\endhead
\bottomrule\noalign{}
\endlastfoot
Beschreibung & \url{https://openalex.org/about} \\
Abdeckung & Aggregatoren sowie Repositorien
(\href{https://openalex.org/sources?page=1&filter=type:repository}{Repositorien-Liste},
\href{https://help.openalex.org/hc/en-us/articles/28932712154391-How-does-OpenAlex-work}{Information
zu Datenquellen},
\href{https://help.openalex.org/hc/en-us/articles/24397285563671-About-the-data}{Noch
eine Information zu Datenquellen}) \\
Metadatenschema &
\href{https://developers.openalex.org/api-reference/introduction}{API
Overview}\_\_Institutionen werden automatisch zugeordnet
(\href{https://help.openalex.org/hc/en-us/articles/24831328396311-Institutions-and-Raw-Affiliation-String-Parsing}{Beschreibung}) \\
Zugänge & \href{https://openalex.org/}{OpenAlex GUI},
\href{https://developers.openalex.org/}{API},
\href{https://developers.openalex.org/download/overview}{Snapshot} ;
Suchfilter Institution über die Nutzeroberfläche (mit Lookup) ; Ab
Februar 2026 ist für die Nutzung der API im real-produktiven Maßstab ein
API-Key erforderlich. Dieser wird nach der Erstellung eines
Nutzeraccounts zur Verfügung gestellt und ist im Basismodell kostenlos.
Das Anfragenlimit beträgt 100 Anfragen pro Sekunde und 100.000 Credits
pro Tag. \\
Suchfelder & Im Suchfilter ``Institution'' können per Lookup
Institutionen ausgewählt werden. Der Publikationstyp (``Type'') kann auf
``dataset'' gefiltert werden. \\
Ausgabeformate & GUI: vollständiger Export nach Login (kostenfrei), csv
oder ``csv excel-optimiert'' (d.~h.
\href{https://help.openalex.org/hc/en-us/articles/24829724234007-Export-results-from-the-OpenAlex-website}{nicht
zu viel Inhalt pro Zelle}) ; REST API: JSON ; Snapshot: JSON, Parquet \\
\end{longtable}
}

\chapter{Verwendung}\label{verwendung}

Einrichtungen verfügen über unterschiedliche technische Möglichkeiten,
die aus den verschiedenen Quellen erhobenen Daten zu verwenden. Je nach
Zielsetzung (s. \hyperref[voruxfcberlegungen]{Kapitel 2}) variieren die
Zwecke und Prioritäten. Dieses Kapitel gibt zunächst einen Überblick
über die möglichen Zielsysteme für die erhobenen Daten und skizziert
anschließend den Ingest\footnote{Mit „Ingest'' ist hier die Einspeisung
  in ein System gemeint. Da der englische Begriff im Deutschen geläufig
  ist, verwenden wir ihn im gesamten Dokument.}-Vorgang mit den zu
berücksichtigenden Entscheidungen und Rahmenbedingungen.

\section{Zielsysteme}\label{zielsysteme}

\subsection{Forschungsinformationssysteme}\label{forschungsinformationssysteme}

Die Verwendung eines FIS als Zielsystem liegt nahe, sofern die eigene
Institution eines betreibt bzw. zur Verfügung stellt.

Es gibt verschiedene Softwarelösungen für ein FIS.\footnote{\url{https://kfid-online.de/fis_landkarte_visual/fis_landkarte_software.php}
  (abgerufen am 30.09.2026)} Im Folgenden wird eine Auswahl von
FIS-Softwarelösungen inkl. einiger Besonderheiten ihres jeweiligen
Umgangs mit Forschungsdaten vorgestellt.

\begin{itemize}
\tightlist
\item
  \href{https://wiki.lyrasis.org/display/DSPACECRIS}{DSpace-CRIS}:
  DSpace-CRIS unterstützt ein erweitertes, an Dublin Core und CERIF
  orientiertes Datenmodell, konfigurierbare Beziehungen, Workflows,
  ORCID-Anbindungen sowie REST-Schnittstellen.\footnote{\url{https://4science.com/english-dspace-cris-7-has-landed/}
    (abgerufen am 30.09.2026)} Ein eigener Dokumenttyptyp für
  Forschungsdatenpublikationen mit lokal individuell konfigurierten
  Metadatenfeldern ist beispielsweise möglich. Mit DSpace 10.0 wurde
  2026 die Integration von DSpace-CRIS-Funktionen in den DSpace-Kern
  begonnen; der Funktionsumfang kann deshalb je nach verwendeter
  Softwareversion deutlich variieren.\footnote{\url{https://dspace.org/dspace-10-0-released/}
    (abgerufen am 30.09.2026)}
\item
  \href{https://www.elsevier.com/solutions/pure}{Pure}: Die Erfassung
  von Forschungsdaten ist als eigener Content-Type „Forschungsdaten'',
  „Datasets'' möglich. Verknüpfungen zu Personen, Organisationseinheiten
  etc. werden in Form von „Master data'' erfasst.\footnote{\url{https://helpcenter.pure.elsevier.com/overview-of-content-types}
    (abgerufen am 30.09.2026)} In Pure lassen sich auch Berichte zum
  Forschungsdatenaufkommen einer Einrichtung umsetzen sowie eine
  (Zweit-)Veröffentlichung der gefundenen Forschungsdatensätze im
  Pure-eigenen Open Science Portal realisieren.
\item
  \href{https://www.symplectic.co.uk/theelementsplatform/}{Elements}:
  Die Erfassung von Forschungsdaten war in Elements seit jeher als
  eigener Datentyp (mit dem Anzeigenamen „Dataset'') möglich. Bei der
  Konfiguration von Elements kann der Datentyp zudem umbenannt werden
  (z. B. in „Forschungsdaten'') oder einen eigenen Datentyp mit einem
  eigenen Satz von Metadatenfeldern erstellen. Für Elements ist
  „Dataset'' ein Untertyp der Kategorie „Publikation''; die Kategorie
  lässt sich jedoch ebenfalls umbenennen (z. B. in „Scholar and Creative
  Work'', falls die semantische Bandbreite erhöht werden soll, bspw. mit
  Blick auf nicht veröffentlichte Forschungsdaten wie solche mit
  „restricted access''). Elements verfügt über eine Integration von
  Figshare, welche die Aufnahme von Metadaten zur Beschreibung von
  Forschungsdaten unterstützt und bei der Aufnahme einen Datensatz
  direkt mit einem Benutzer im Elements-Objektmodell verknüpft. Ebenso
  kann Elements so konfiguriert werden, dass es Metadaten zur
  Beschreibung von Forschungsdaten aus dem institutionellen Repositorium
  einer Einrichtung übernimmt. Derzeit unterstützt Elements
  Integrationen von DSpace, EPrints, Hyrax und Figshare for
  Institutions. Mittels der Elements-Benutzeroberfläche lassen sich
  Metadaten zur Beschreibung von Forschungsdaten hinzufügen; alternativ
  können sie über die Elements-API aus anderen strukturierten
  Datenquellen bezogen werden, sofern entsprechende Zugriffsrechte
  vorliegen.
\item
  \href{https://clarivate.com/academia-government/ko/scientific-and-academic-research/research-funding-analytics/converis/}{Converis}:
  In der FIS-Software sind ebenfalls schon seit jeher individuelle
  Anpassungen am Datenmodell und Metadatenschema möglich; das System
  kann Publikationen, Personen, Organisationseinheiten, Projekte und
  weitere Forschungsinformationen zusammenführen und mit Repositorien
  sowie öffentlichen Forschungsportalen verbinden. Die Converis-Instanz
  der Universität Kassel enthält bspw. den Publikationstyp
  „Forschungsdatenpublikation''.\footnote{\url{https://forschung.uni-kassel.de/converis/portal/list?show=Publication&filter=\%5Bfpublyear\%3D\%5D\%2C\%5Bfpublicationtype\%3D15998\%5D&sortBy=publYear&sortOrder=false&cypher=All&page=1&items=10&showAll=&treeView=false&moreOrLess=\%5Bfpublyear\%3Dtrue\%5D\%2C\%5Bfpublicationtype\%3Dfalse\%5D&auxfun=&lang=de_DE}
    (abgerufen am 30.09.2026)} Das FIS der
  Friedrich-Alexander-Universität Erlangen-Nürnberg basiert ebenfalls
  auf Converis und erfasst Forschungsdaten; die Entwickler:innen können
  zum Erfahrungsaustausch angesprochen werden.
\item
  \href{https://www.his.de/hisinone/forschungsmanagement}{HISinOne-RES}:
  Das System enthält Funktionalitäten für ein umfassendes
  Forschungsmanagement und orientiert sich beim Metadatenschema eng am
  Kerndatensatz Forschung (KDSF). Der KDSF-Systematik entsprechend gibt
  es einen Publikationstyp „Forschungsdaten'', der mit anderen Objekten
  wie Personen, Projekten und Förderinformationen verknüpft werden kann.
  Außerdem bietet die Software Optionen des automatisierten Ingests per
  API. Importmöglichkeiten bestehen via ORCID, RIS, BibTeX, OpenAlex und
  OAI-PMH. Um den Import von Metadaten aus Radar Local (RADAR XML) per
  API zu ermöglichen, hat die Universitätsbibliothek Würzburg ein
  XSLT-Stylesheet zur Metadatentransformation erstellt, welches seit
  Januar 2026 zusammen mit dem FIZ Karlsruhe auch zur Nachnutzung
  bereitgestellt wird.\footnote{\url{https://github.com/ER-FIZKarlsruhe/radar-stylesheets}
    (abgerufen am 30.09.2026)}
\end{itemize}

\subsection{Institutionelles
Repositorium}\label{institutionelles-repositorium}

Forschungsdaten bzw. ihre Metadaten können auch im institutionellen
Repositorium, das entweder ein reines Forschungsdaten- oder ein
gemischtes Daten- und Textrepositorium sein kann, verzeichnet werden.
Institutionen entscheiden sich für das institutionelle Repositorium als
Zielsystem beispielsweise, weil die Einrichtung kein
Forschungsinformationssystem nutzt oder weil eine regelmäßige
automatisierte Übertragung von Publikationsmetadaten vom Repositorium
ins Forschungsinformationssystem bereits etabliert ist. Im zweiten Fall
kann die bestehende Integration für die aus externen Quellen
geharvesteten Metadaten ohne weitere Entwicklungsarbeit auf
Forschungsdaten ausgeweitet werden.

Beispiele für Forschungsdatenrepositorien als Zielsysteme finden sich
bei:

\begin{itemize}
\tightlist
\item
  \href{https://dataverse.lib.virginia.edu/dataverse/libradata?q=&fq2=metadataSource\%3A\%5C\%22Harvested\%5C\%22&fq0=subtreePaths\%3A\%5C\%22\%2F1\%5C\%22&fq1=dvObjectType\%3A\%28dataverses+OR+datasets\%29&types=dataverses\%3Adatasets&sort=dateSort&order=}{University
  of Virginia Dataverse}
\item
  \href{https://dataverse.harvard.edu/dataverse/harvard;jsessionid=3945204ee9c13b5ef80975f097f4?q=&fq1=metadataSource\%3A\%5C\%22Harvested\%5C\%22&fq0=dvObjectType\%3A\%28dataverses+OR+datasets\%29&types=dataverses\%3Adatasets&sort=dateSort&order=}{Harvard
  Dataverse}
\item
  \href{https://university.open.ac.uk/library-research-support/research-data-management/open-university-research-data-management-policy}{The
  Open University}: „4.2 Where data have been archived in a repository
  other than ORDO or the OU Scholarship Exchange, a metadata-only record
  should be published on ORDO which describes the data and provides a
  persistent link (such as a DOI) to the data.''
\item
  Baldwin und Green (2024)
\item
  \href{https://www.uef.fi/en/open-research-data}{University of Eastern
  Finland}
\end{itemize}

Die folgenden Softwarelösungen für Repositorien sind derzeit weit
verbreitet:

\begin{itemize}
\tightlist
\item
  \href{https://www.mycore.de/en/}{MyCoRe}
\item
  \href{https://dspace.lyrasis.org/}{DSpace} (das auch eine FIS-Version
  anbietet, siehe oben). In DSpace können diverse Publikationstypen mit
  jeweils eigenem Metadatensatz konfiguriert werden, ab Version 9 auch
  als eigenständige Entitäten, die miteinander in Beziehung gesetzt
  werden können.
\item
  \href{https://www.eprints.org/uk/}{EPrints} (EPrints (2015))
\item
  \href{https://inveniosoftware.org/}{Invenio} (als
  Open-Source-Framework für groß anzulegende digitale Repositorien oder
  spezifisch InvenioRDM für Forschungsdaten)
\item
  \href{https://www.fiz-karlsruhe.de/de/produkte-und-dienstleistungen/radar}{RADAR}
  (In den Betriebsvarianten RADAR Local und RADAR Cloud)
\end{itemize}

Gegebenenfalls muss vor der Erfassung von Forschungsdaten zunächst ein
Dokumenttyp oder eine Entität für Forschungsdaten angelegt werden.

\subsection{Hochschulbibliografie und
Bibliothekskatalog}\label{hochschulbibliografie-und-bibliothekskatalog}

Des Weiteren können Forschungsdaten in der Hochschulbibliografie
verzeichnet oder in den Bibliothekskatalog aufgenommen werden.

Die Universität Heidelberg verzeichnet beispielsweise in ihrer
Hochschulbibliografie heiBIB auch
Forschungsdatenpublikationen.\footnote{\url{https://heibib.ub.uni-heidelberg.de/search/Search/Results?type=AllFields&filter\%5B\%5D=format\%3A\%22Research+Data\%22}
  (abgerufen am 30.09.2026)} Ein weiteres Beispiel ist die Universität
Bamberg, die seit Anfang 2022 Forschungsinformationen zu Forschungsdaten
in ihrer Universitätsbibliografie im FIS verzeichnet und mit
Forschungsinformationen zu Personen, Einrichtungen, Projekten,
Publikationen und Auszeichnungen verknüpft (Ziegler und Franke (2023) ;
S. 1 - 13).

Im gemeinsamen Katalog der Bibliotheksverbünde BVB und KOBV, B3Kat,
können Forschungsdatenpublikationen manuell erfasst werden.\footnote{\url{https://www.bib-bvb.de/web/kkb-online/forschungsdaten}
  (abgerufen am 30.09.2026)} Das Discovery-System EBSCO, das Daten an
mehrere Bibliothekskataloge liefert, bezieht Metadaten unter anderem von
der Bielefeld Academic Search Engine (BASE), die auch
Forschungsdatenpublikationen verzeichnet. Auf diese Art gelangen
Datenpublikations-Metadaten also auch automatisch über Umwege in manche
Bibliothekskataloge.

Eine Art Hochschulbibliografie für Daten führt die University of
Helsinki mit ihrem Data catalogue. „The metadata are gathered in Data
catalogue through automatic data transfer (harvesting), and user
submissions. For harvestable repositories, the researcher must have the
"University of Helsinki" in their affiliation for the metadata to appear
in the Data catalogue.''\footnote{\url{https://www.helsinki.fi/en/research/services-researchers/data-support/preservation-research-data/data-catalogue}
  (abgerufen am 30.09.2026)}

\section{Schritte zum Ingest in das
Zielsystem}\label{schritte-zum-ingest-in-das-zielsystem}

Nach den Vorarbeiten (s. \hyperref[voruxfcberlegungen]{Kapitel 2}) und
der Datensuche (s. \hyperref[durchfuxfchrung]{Kapitel 3}) sollen die
identifizierten Publikationsmetadaten nun in das anvisierte Zielsystem
überführt werden. Unabhängig vom konkreten Zielsystem umfasst der
Ingest-Prozess die folgenden Schritte.

\begin{itemize}
\tightlist
\item
  \textbf{Abbildung, Sichtbarkeit und Berechtigungen festlegen} Um die
  erfassten Forschungsdaten importieren oder händisch eintragen zu
  können, muss sichergestellt sein, dass passende Metadatenfelder im
  Zielsystem existieren oder geschaffen werden können, idealerweise auf
  Basis bestehender gängiger Standards wie dem Kerndatensatz
  Forschung.\footnote{\url{https://www.kerndatensatz-forschung.de/kdsf20/kdsf20.php}
    (abgerufen am 30.09.2026)} Das Zielsystem hat eine bestimmte
  Struktur mit festgelegten Elementen. Ein CRIS kann beispielsweise in
  Sammlungen für Fachbereiche organisiert sein und Entitäten wie
  „Artikel'', „Dissertation'', „Person'', „Projekt'' u. a. enthalten.
  Wie sollen die Metadaten der Datenpublikationen abgebildet werden?
  Gibt es bspw. einen Elementtyp „Dataset'' oder „Datensatz''? Die
  gewünschte Darstellungsform ist bei den weiteren Überlegungen
  (Ingest-Optionen, Metadatenschema) zu berücksichtigen. Ebenso sollte
  vorab geplant werden, welche Beschränkungen und Berechtigungen für die
  eingespielten Daten gelten sollen (read / write). Je nachdem, wie
  diese im Zielsystem organisiert sind, können oder sollten die
  Berechtigungen direkt beim Ingest gesetzt werden, bspw. über
  entsprechende Metadatenfelder oder zusätzliche Angaben im
  Ingest-Prozess.
\item
  \textbf{Ingest-Optionen des Zielsystems prüfen} In den meisten Fällen
  ist eine manuelle Eingabe der Publikationsmetadaten nicht
  wünschenswert. In manche Systeme (z. B. DSpace) können Daten in
  strukturierter Form (z. B. als CSV-Datei mit bestimmten
  Spaltenbezeichnungen, oder im OAIS-Format) importiert werden.
  Möglicherweise steht auch eine API für den Import zur Verfügung. Für
  den Import müssen ggf. die nötigen Berechtigungen erteilt werden.
\item
  \textbf{Metadatenschema des Zielsystems prüfen} Das Metadatenschema
  des Zielsystems kann von denen der genutzten Datenquellen abweichen.
  Beispielsweise können die Namen der Autor:innen im Metadatenfeld
  dc.contributor.author stehen, in dcterms.creator oder in Feldern aus
  noch anderen Metadatenschemata. Damit die Metadaten korrekt in das
  Zielsystem übertragen werden, muss ein Mapping zwischen dem
  Metadatenschema der Treffermenge auf das des Zielsystems erstellt
  werden.
\item
  \textbf{Datensatz-Metadaten auf Zielsystem anpassen} Anschließend
  werden die Publikationsmetadaten transformiert, um a) dem
  Metadatenschema des Zielsystems und b) ggf. der für den Ingest
  erforderlichen strukturierten Form zu entsprechen. Für diesen Schritt
  ist eine gute Dokumentation von besonderer Bedeutung, um eine
  wiederholte oder regelmäßige Durchführung vorzubereiten. Die Dokumente
  für die Transformation können ggf. auch veröffentlicht werden, um die
  Nachnutzung durch Dritte zu ermöglichen.
\item
  \textbf{Ingest testen} Vor dem Komplett-Ingest ist ein Test mit
  begrenztem Umfang (z. B. 10 Publikationen) unbedingt empfehlenswert.
\item
  \textbf{Ingest durchführen} War der Test erfolgreich, kann der Ingest
  vollständig durchgeführt und anschließend geprüft werden.
\item
  \textbf{Überprüfen und dokumentieren} Eine Überprüfung des Ingests,
  entweder stichprobenartig oder automatisiert und vollständig, ist
  angeraten. Eine gute Dokumentation hilft, wie immer, sowohl bei der
  Identifizierung von Fehlerquellen als auch bei einer potentiellen
  erneuten Ausführung zu einem späteren Zeitpunkt.
\end{itemize}

\chapter{Vorbereitung}\label{vorbereitung}

Bestimmte Rahmenbedingungen können die Erfassung des
Forschungsdatenoutputs an einer Einrichtung erleichtern. Das ist
insbesondere bei permanentem Monitoring relevant, also wenn die
Erfassung in regelmäßigen Abständen wiederholt werden soll. Im Folgenden
wird auf einige dieser Rahmenbedingungen eingegangen und beschreiben
Maßnahmen beschrieben, die zukünftige Monitoringaktivitäten unterstützen
können. Die beschriebenen Maßnahmen betreffen neben dem Sammeln der
Metadaten auch die Motivation und Befähigung von Forschenden.

\section{Organisatorische Maßnahmen für
Sammlungen}\label{organisatorische-mauxdfnahmen-fuxfcr-sammlungen}

Mit gezielten Maßnahmen können Prozesse und Infrastrukturen besser auf
die Erfassung von Forschungsdatenpublikationen vorbereitet werden. Da
Voraussetzungen an jeder Einrichtung sehr vielseitig sein können, dienen
die genannten Beispiele primär der Orientierung.

\subsection{Leit- und Richtlinien}\label{leit--und-richtlinien}

Um einen Rahmen für das Monitoring von Forschungsdatenpublikationen zu
schaffen und entsprechende Abfragen zu vereinfachen, können
Einrichtungen die Absicht einer möglichst vollständigen Erfassung in
Open-Science- oder Forschungsdatenleitlinien aufnehmen. Eine
verpflichtende Erfassung sieht beispielsweise die Research Data
Management Policy der Buckinghamshire New University vor: „All research
data must be registered with the University whether it is hosted by the
University or retained elsewhere'' (Buckinghamshire New University
(2016) ; S. 3). Hierbei lassen sich auch Regelungen zur Nutzung des
institutionellen Repositoriums und zu Communities der Institution in
externen Repositorien wie Zenodo festhalten.

Formulierte Vorgaben und Handreichungen zielen darauf ab, eine
bestmöglichen Angabe von Metadaten bei Forschungsdatenpublikationen
durch die Autor:innen der eigenen Einrichtung zu erreichen. Dies umfasst
neben der korrekten Affiliationsangabe, etwa festgelegt durch eine
Affiliationsrichtlinie, auch die Förderung der Nutzung von persistenten
Identifikatoren für Einrichtungen (ROR) und Autor:innen (ORCID).
Affiliationsrichtlinien werden in der Regel nach Abstimmung mit dem
Präsidium erarbeitet und vom Senat verabschiedet. Viele universitäre
Einrichtungen haben eine Affiliationsrichtlinie, beispielsweise die TU
Clausthal.\footnote{\url{https://doi.org/10.21268/20260119-0}}

\subsection{Institutsseitige Empfehlungen für die
Repositorienwahl}\label{institutsseitige-empfehlungen-fuxfcr-die-repositorienwahl}

Es gibt eine Vielzahl von Forschungsdatenrepositorien, die sich mit
Blick auf ihren Anwendungsbereich, die technische Basis, Implementierung
von Standards und weiteren Eigenschaften unterscheiden. Das Verzeichnis
re3data weist inzwischen mehr als 3500 Repositorien nach.\footnote{\href{https://www.re3data.org/}{https://www.re3data.org}
  (abgerufen am 30.09.2026)} Forschungsförderer und Akteure wie z. B.
DINI definieren vermehrt Anforderungen an geeignete
Forschungsdatenrepositorien, um Forschende Orientierung bei der
Repositorienwahl zu geben (Gollwitzer u.~a. (2020)). Auch Einrichtungen
können Empfehlungen aussprechen, wie etwa die ETH Zürich in den
„Guidelines for Research Data Management at ETH Zurich''\footnote{\url{https://ethz.ch/content/dam/ethz/main/eth-zurich/organisation/rechtssammlung/414.2en.pdf}
  (abgerufen am 30.09.2026)} oder die Forschungsdatenleitlinie der
Max-Weber-Stiftung\footnote{\url{https://doi.org/10.5281/zenodo.10219502}}.

Durch Empfehlungen wie diese können Einrichtungen ihre Angehörigen an
Forschungsdatenrepositorien verweisen, die gewünschte Eigenschaften
aufweisen. Hierzu zählen beispielsweise die Nutzung persistenter
Identifikatoren oder das Harvesten der Datensätze durch Aggregatoren.
Wenn Forschende ihre Daten dort veröffentlichen, wird die zukünftige
Erfassung des Publikationsoutputs erleichtert.

\subsection{Technische Voraussetzungen
schaffen}\label{technische-voraussetzungen-schaffen}

Je nach gewünschtem Verfahren (s. \hyperref[durchfuxfchrung]{Kapitel 3})
und Voraussetzungen des Zielsystems (s. \hyperref[verwendung]{Kapitel
4}) können (semi-)automatisierte Importroutinen implementiert werden.
Das für die Erfassung zuständige Team kann unabhängig von Forschenden
Einträge importieren. Auch manuelle Such- und Erfassungsvorgänge sind --
mit entsprechendem Aufwand -- denkbar, etwa für einen anlassbezogenen
Überblick (s. \hyperref[ziele]{Kapitel 2.1}).

\subsection{Teilnahme am ORCID-DE
Konsortium}\label{teilnahme-am-orcid-de-konsortium}

Sofern es keine vertragliche Verpflichtung gibt, obliegt die
Entscheidung, einen persönlichen ORCID-Account einzurichten und zu
pflegen, jedem Forschenden selbst. Während das Einrichten und die Angabe
der ORCID meist ohne Unterstützungs- oder Korrekturbedarf abläuft,
erfolgt die aktive Pflege der Einträge nicht immer konsequent.
Einrichtungen, die Mitglied im ORCID-DE Konsortium sind, erhalten die
Möglichkeit, die ORCID-Einträge ihrer Forschenden zu kuratieren und sie
somit bei der Pflege des Accounts zu unterstützen. Der Zugriff auf die
ORCID-Einträge erfolgt von technischer Seite über eine API und bedarf
der expliziten Zustimmung des jeweiligen Forschenden. Weitere
Informationen zum Vorgang sind unter
\url{https://www.orcid-de.org/konsortium/ueber-das-konsortium} und auf
den verknüpften Seiten zu finden.

Durch diese Serviceleistung kann die Akzeptanz der ORCID bei Forschenden
erhöht und zugleich die Metadatenqualität in den institutseigenen
Forschungsinformationssystemen verbessert werden. Dennoch ist an dieser
Stelle anzumerken, dass die Konsortialmitgliedschaft nicht kostenlos
ist. Einrichtungen müssen hier abwägen, ob sie den Mitgliedsbetrag
aufbringen können bzw. wollen. Auch der erforderliche Personaleinsatz
für die Datenpflege sollte bedacht werden.

\section{Motivation und Befähigung von
Forschenden}\label{motivation-und-befuxe4higung-von-forschenden}

Die Einbindung von Forschenden in den Prozess der Sammlung und
Sichtbarmachung des Forschungsoutputs der eigenen Einrichtung bereits
während des Publikationsprozesses kann zu einer Verbesserung der
Metadatenqualität führen und so das spätere Auffinden der Datensätze
sowie das Transferieren in andere Systeme erleichtern. Wesentliche
Punkte bei diesem Vorgehen sind die Motivation und die Befähigung der
Forschenden, sich aktiv in den Prozess einzubringen. Hierfür müssen
seitens der Einrichtung die Rahmenbedingungen in Form von Leit- und
Richtlinien ausgearbeitet, verabschiedet und veröffentlicht werden.
Außerdem können Anreize geschaffen werden, die das Mitwirken der
Forschenden positiv beeinflussen. Um die Forschenden über ihre Rolle im
Prozess, das damit verbundene Ziel und die zugrundeliegenden Bedingungen
aufzuklären sowie dafür zu sensibilisieren, ist eine zielgerichtete
Informationsbereitstellung von großer Bedeutung.

\subsection{Anreizsysteme}\label{anreizsysteme}

Die Publikation von Forschungsdaten sowie die damit verknüpfte Beachtung
der Leit- und Richtlinien lassen sich durch bestimmte Anreizmechanismen
fördern.

Durch die Berücksichtigung von Open-Science-Praktiken in Einstellungs-
und Evaluationsverfahren oder in der leistungsorientierten Mittelvergabe
können Anreize für die Veröffentlichung von Forschungsdaten gesetzt
werden. Dies erfolgt bereits an einigen Einrichtungen, wie eine
internationale Sammlung von Stellenausschreibungen zeigt (Schönbrodt
u.~a. (2018)). Darüber hinaus können Datenpublikationen inkl. korrekter
Affiliationsangaben durch Preise honoriert werden.

\subsection{Informationsangebote für
Forschende}\label{informationsangebote-fuxfcr-forschende}

Um die Forschenden als aktive Akteure über die anvisierte Sammlung und
die damit designierte Verbesserung der Sichtbarkeit von Forschungsdaten
der eigenen Einrichtung zu sensibilisieren, sind zielgerichtete
Informationen -- auch über die am jeweiligen Standort herrschende
Verbindlichkeit der Umsetzung - essentiell.

Die Informationsangebote können die folgenden Aspekte in Bezug auf die
Datenpublikation behandeln (soweit zutreffend):

\begin{itemize}
\tightlist
\item
  die Nutzung persistenter Identifier (insbesondere DOI, ORCID, ROR)
\item
  die Wahl eines Repositoriums (Auswahlkriterien, empfohlene
  Repositorien, Verweis auf das institutionelle Repositorium)
\item
  Institutsseitige Vorgaben

  \begin{enumerate}
  \def\labelenumi{\arabic{enumi}.}
  \tightlist
  \item
    Forschungsdatenleitlinie oder Open-Science-Leitlinie
  \item
    Affiliationsrichtlinie
  \item
    Leitlinien für die Datenaufnahme und -kuration in Repositorien (z.
    B. in Zenodo-Communities)
  \item
    Vorgaben zur Dateneinspeisung ins FIS
  \item
    Honorierung guter Praxis bei Einstellungsverfahren
  \end{enumerate}
\item
  FAIR Data Principles
\end{itemize}

Forschende können hierbei über unterschiedliche Kommunikationskanäle
erreicht werden. Dafür bieten sich die folgenden Wege an:

\begin{itemize}
\tightlist
\item
  Internetseiten des lokalen Forschungsdatenmanagements oder der
  zugehörigen Universitätsbibliothek
\item
  Unterlagen in Onboarding- und Offboarding-Verfahren
\item
  Newsletter
\item
  Leitlinien und Handlungsempfehlungen
\end{itemize}

Eine dynamische Informationsbereitstellung kann unter anderem durch die
direkte Interaktion mit den Forschenden auf Informationsveranstaltungen,
Multiplikatorveranstaltungen (z. B. Sitzungstermine im Senat oder
Fakultätsrat), in Onboarding-Gesprächen und in gezielten Schulungen
erfolgen.

Hierbei ist hervorzuheben, dass erfahrungsgemäß insbesondere
(Post-)Doktoranden gut über Schulungen und Workshops zu Themen des
Forschungsdatenmanagements erreichbar sind, aber auch andere Zielgruppen
können durch spezialisierte Angebote angesprochen werden.

Neben dem FDM- und Bibliothekspersonal werden hinsichtlich der
Kommunikation mit den Forschenden aber auch die Präsidiumsmitglieder,
welche für das Ressort Forschung zuständig sind, in der Verantwortung
gesehen, sich unterstützend einzubringen. Eine vom Präsidiumsmitglied
verfasste Rundmail, in der die Forschenden dazu ermutigend aufgefordert
werden, eine ORCID anzulegen und zu pflegen, die Affiliationsrichtlinie
zu beachten oder bestimmte Vorgaben bei der Datenpublikation zu
erfüllen, kann eine verstärkende Wirkung besitzen.

\chapter{Zusammenfassung}\label{zusammenfassung}

Abschließend lässt sich festhalten, dass ein institutionelles Monitoring
dezentral publizierter Forschungsdaten realisierbar, aber mit einigen
Herausforderungen verbunden ist.

Zu den Herausforderungen gehört insbesondere der Ressourceneinsatz, der
für die Umsetzung erforderlich ist, und die Limitationen der verwendeten
Datenquellen, z. B. mit Blick auf deren Datenqualität oder
Veränderlichkeit (Demes (2026)). Auch das Policy-Umfeld ist geprägt von
dynamischen Entwicklungen. So diskutieren verschiedene Akteure die
Berücksichtigung von Open-Science-Praktiken in Evaluations- und
Berufungsverfahren und setzen damit Anreize für die Erfassung des
Forschungsdatenoutputs, und die internationale Open Science Monitoring
Initiative fördert den Austausch und setzt sich für mehr
Standardisierung ein. Diese Entwicklungen können schwer zu überblicken
sein.

Das Monitoring des institutionellen Forschungsdatenoutputs ist ein
umfangreiches Vorhaben, das nicht nur ein Verständnis für
institutionelle Gegebenheiten, Datenquellen und Suchstrategien\st{,}
sondern auch zumindest grundlegende technische Fähigkeiten voraussetzt.
Das Monitoring betrifft darüber hinaus verschiedene Bereiche der
Einrichtungen und profitiert darum von klaren Zuständigkeiten und
Vorgaben.

Diese Handreichung zeigt auf, dass erste Schritte möglich sind, die
schrittweise ausgebaut und so immer bessere Ergebnisse liefern können.
Das Dokument sammelt Erfahrungen und bereitet sie in einer Form auf, die
hoffentlich anderen Interessierten den Einstieg in das Thema
erleichtert. Dieses Dokument wird zukünftig ergänzt um Notebooks, die
konkrete Beispiele für Suchanfragen liefern.

\chapter*{Literaturverzeichnis}\label{literaturverzeichnis}
\addcontentsline{toc}{chapter}{Literaturverzeichnis}

\markboth{Literaturverzeichnis}{Literaturverzeichnis}

\protect\phantomsection\label{refs}
\begin{CSLReferences}{1}{1}
\bibitem[\citeproctext]{ref-baldwin_instant_2024}
Baldwin, Julie, und Jonathan Green. 2024. \emph{Instant Data Finding and
cataloguing externally hosted research datasets through automation}.
Figshare. \url{https://doi.org/10.6084/M9.FIGSHARE.25233181}.

\bibitem[\citeproctext]{ref-bobrov_principles_2025}
Bobrov, Evgeny, Laetitia Bracco, Marin Dacos, u.~a. 2025. \emph{The
Principles of Open Science Monitoring}.
\url{https://doi.org/10.5281/zenodo.17853227}.

\bibitem[\citeproctext]{ref-buckinghamshire_new_university_research_2016}
Buckinghamshire New University. 2016. \emph{Research Data Management
Policy}.
\url{https://www.bucks.ac.uk/sites/default/files/2021-03/research_data_management_policy.pdf}.

\bibitem[\citeproctext]{ref-cohen_matching_2026}
Cohen, Avihay, Blanka Ivanovic, Anastasiia Iarkaeva, Vladislav Nachev,
und Evgeny Bobrov. 2026. \emph{Matching data references and
institutional output to map the reuse of biomedical research data}.
z9kjf\_v4. {MetaArXiv}. \url{https://doi.org/10.31222/osf.io/z9kjf_v4}.

\bibitem[\citeproctext]{ref-datacite_datacite_2026}
DataCite. 2026. \emph{{DataCite} Public Data File 2025}. {DataCite}.
\url{https://doi.org/10.14454/T5QB-D995}.

\bibitem[\citeproctext]{ref-datacite_data_2025}
DataCite, und Make Data Count. 2025. \emph{Data Citation Corpus Data
File}. Version v3.0. Zenodo.
\url{https://doi.org/10.5281/zenodo.14897662}.

\bibitem[\citeproctext]{ref-demes_overhaul_2026}
Demes, Kyle. 2026. {„An Overhaul of Type Classification. {OpenAlex}
blog``}. Juli 15.
\url{https://blog.openalex.org/an-overhaul-of-type-classification/}.

\bibitem[\citeproctext]{ref-eprints_eprints_2015}
EPrints. 2015. {„{EPrints} for Research Data-- {EPrints} Services``}.
Juli 9.
\url{https://uk.eprints-hosting.org/uk/index.php/flavours/eprints-for-research-data-1/}.

\bibitem[\citeproctext]{ref-gee_hunt_2026}
Gee, Bryan M. 2026. {„The hunt for research data: Development of an
open-source workflow for tracking institutionally-affiliated research
data publications``}. \emph{Journal of {eScience} Librarianship} 15 (1).
\url{https://doi.org/10.7191/jeslib.1170}.

\bibitem[\citeproctext]{ref-gerasimov_comparison_2024}
Gerasimov, Irina, Binita KC, Armin Mehrabian, James Acker, und Michael
P. McGuire. 2024. {„Comparison of datasets citation coverage in Google
Scholar, Web of Science, Scopus, Crossref, and {DataCite}``}.
\emph{Scientometrics} 129 (7): 3681--704.
\url{https://doi.org/10.1007/s11192-024-05073-5}.

\bibitem[\citeproctext]{ref-gollwitzer_management_2020}
Gollwitzer, Mario, Andrea Abele-Brehm, Christian Fiebach, u.~a. 2020.
\emph{Management und Bereitstellung von Forschungsdaten in der
Psychologie: Überarbeitung der {DGPs}-Empfehlungen}. {PsyArXiv}.
\url{https://doi.org/10.31234/osf.io/hcxtm}.

\bibitem[\citeproctext]{ref-gregory_tracing_2023}
Gregory, Kathleen, Anton Ninkov, Chantal Ripp, Emma Roblin, Isabella
Peters, und Stefanie Haustein. 2023. {„Tracing data: A survey
investigating disciplinary differences in data citation``}.
\emph{Quantitative Science Studies} 4 (3): 622--49.
\url{https://doi.org/10.1162/qss_a_00264}.

\bibitem[\citeproctext]{ref-habermann_datacite_2025}
Habermann, Ted. 2025. {„{DataCite} Facets and Metadata Completeness.
Metadata Game Changers``}. \url{https://doi.org/10.59350/55etd-e8154}.

\bibitem[\citeproctext]{ref-iarkaeva_methodology_nodate}
Iarkaeva, Anastasiia, Maaike Duine, und Andreas Hübner. o.~J.
\emph{Methodology to detect Open Data underlying the journal articles'
full texts of the Department of Earth Sciences, Freie Universität
Berlin}. \url{https://doi.org/10.57874/Z5FK-WE46}.

\bibitem[\citeproctext]{ref-iarkaeva_workflow_2024}
Iarkaeva, Anastasiia, Vladislav Nachev, und Evgeny Bobrov. 2024.
{„Workflow for detecting biomedical articles with underlying open and
restricted-access datasets``}. \emph{{PLOS} {ONE}} 19 (5): e0302787.
\url{https://doi.org/10.1371/journal.pone.0302787}.

\bibitem[\citeproctext]{ref-kirsch_assessing_2026}
Kirsch, Danielle R., und Isaac Wink. 2026. {„Assessing Researcher Data
Sharing Practices Using Publicly Available Dataset Metadata: A
Reproducible Workflow for Institutional Stakeholders``}. \emph{Journal
of Librarianship and Scholarly Communication} 14 (1).
\url{https://doi.org/10.31274/jlsc.22913}.

\bibitem[\citeproctext]{ref-putnings_go_2024}
Putnings, Markus, Anna Daudrich, Nils Dille, Anastasiia Iarkaeva, Janina
Kühner, und Alica Wollmann. 2024. \emph{{GO} {UNITE}! Empfehlungen zur
Sammlung und Sichtbarkeit von Forschungsdaten der eigenen Institution}.
Zenodo. \url{https://doi.org/10.5281/zenodo.11102368}.

\bibitem[\citeproctext]{ref-riedel_oddpub_2026}
Riedel, Nico, und Vladislav Nachev. 2026. \emph{{ODDPub}: A text-mining
algorithm to detect data sharing in biomedical publications}. Version
7.3.0. Released. \url{https://doi.org/10.5281/zenodo.3741403}.

\bibitem[\citeproctext]{ref-robinson-garcia_datacite_2017}
Robinson-Garcia, Nicolas, Philippe Mongeon, Wei Jeng, und Rodrigo
Costas. 2017. {„{DataCite} as a novel bibliometric source: Coverage,
strengths and limitations``}. \emph{Journal of Informetrics} 11 (3):
841--54. \url{https://doi.org/10.1016/j.joi.2017.07.003}.

\bibitem[\citeproctext]{ref-van_de_sandt_definition_2019}
Sandt, Stephanie van de, Sünje Dallmeier-Tiessen, Artemis Lavasa, und
Vivien Petras. 2019. {„The Definition of Reuse``}. \emph{Data Science
Journal} 18 (1): 22. \url{https://doi.org/10.5334/dsj-2019-022}.

\bibitem[\citeproctext]{ref-schonbrodt_academic_2018}
Schönbrodt, Felix, Leonhard Schramm, Franka Etzel, u.~a. 2018.
\emph{Academic job offers that mentioned open science}. {OSF}.
\url{https://doi.org/10.17605/OSF.IO/7JBNT}.

\bibitem[\citeproctext]{ref-strecker_how_2025}
Strecker, Dorothea, Kerstin Soltau, und Felix Bach. 2025. {„How Are
Research Data Referenced? The Use Case of the Research Data Repository
{RADAR}``}. \emph{Research Data Management: Challenges in a Changing
World} (Heidelberg), 146--59.
\url{https://doi.org/10.11588/heibooks.1652}.

\bibitem[\citeproctext]{ref-noauthor_unesco_2021}
\emph{{UNESCO} Recommendation on Open Science}. 2021. {UNESCO}.
\url{https://doi.org/10.54677/MNMH8546}.

\bibitem[\citeproctext]{ref-noauthor_ub-mannheimmadabi_2026}
Universitätsbibliothek Mannheim. 2026. \emph{{UB}-Mannheim/madabi}.
Released. \url{https://github.com/UB-Mannheim/madabi}.

\bibitem[\citeproctext]{ref-ziegler_6_2023}
Ziegler, Barbara, und Fabian Franke. 2023. {„6 in 1: Publikationen,
Forschungsdaten, Projekte, Forschende, Einrichtungen, Auszeichnungen``}.
\emph{o-bib. Das offene Bibliotheksjournal} 10 (1): 1--13.
\url{https://doi.org/10.5282/O-BIB/5868}.

\end{CSLReferences}




\end{document}
